\item \textbf{Colisión elástica de dos bolas de acero.}
Dos bolas de acero idénticas, cada una con una masa $m = 67.4\text{ g}$ y un diámetro $d = 25.4\text{ mm}$, se mueven en direcciones opuestas, cada una a $v = 5.00\text{ m/s}$. Chocan de frente y se separan elásticamente.
\begin{enumerate}
    \item Divida a su grupo en dos y haga que cada mitad encuentre el intervalo de tiempo total en el que las bolas están en contacto, usando diferentes modelos:
    \begin{itemize}
        \item \textbf{Grupo (i):} modele una pelota dada como poseedor de energía cinética que luego se transforma por completo en energía potencial elástica en el instante en que las bolas se detienen momentáneamente. Suponga que la aceleración de la bola durante este intervalo de tiempo es constante y use el modelo de aceleración constante para encontrar el intervalo de tiempo total en el que las bolas están en contacto.
        \item \textbf{Grupo (ii):} apretando una de las bolas en una prensa mientras se hacen mediciones precisas de la cantidad resultante de compresión, se ha descubierto que la ley de Hooke es un buen modelo del comportamiento elástico de la pelota. Una fuerza de $F = 16.0\text{ kN}$ ejercida por cada mordaza del tornillo de banco reduce el diámetro en una distancia $s = 0.200\text{ mm}$. El diámetro vuelve al valor original cuando se elimina la fuerza. Modele el movimiento de cada bola, mientras las bolas están en contacto, como la mitad de un ciclo de movimiento armónico simple. Encuentre el intervalo de tiempo total en el que las bolas están en contacto.
    \end{itemize}
    \item ¿Qué resultado es más exacto?
\end{enumerate}